\documentclass[twocolumn]{revtex4-2}

\usepackage{amsmath,amssymb}
\usepackage{graphicx}
\usepackage{booktabs}
\usepackage{xcolor}
\usepackage{hyperref}

\newcommand{\Nf}{N_f}
\newcommand{\Csq}{C^{2}}
\newcommand{\GB}{\mathcal{G}}
\newcommand{\Riem}{R_{\mu\nu\rho\sigma} R^{\mu\nu\rho\sigma}}
\newcommand{\Ric}{R_{\mu\nu} R^{\mu\nu}}
\newcommand{\Rsq}{R^{2}}

% Auto-generated by update_counts.py — 2026-04-24T19:44:03.140400
% Do not edit manually. Run: uv run python scripts/update_counts.py
\newcommand{\totaltheorems}{3993}
\newcommand{\substantivetheorems}{3882}
\newcommand{\placeholdertheorems}{111}
\newcommand{\axiomcount}{1}
\newcommand{\sorrycount}{0}
\newcommand{\sorrypercent}{0.0}
\newcommand{\leanmodules}{167}
\newcommand{\leandefinitions}{3297}
\newcommand{\pythonmodules}{68}
\newcommand{\testfiles}{59}
\newcommand{\totaltests}{2836}
\newcommand{\figurecount}{111}
\newcommand{\notebookcount}{54}
\newcommand{\papercount}{19}
\newcommand{\aristotleproved}{322}
\newcommand{\aristotleruns}{44}
% --- Paper 18 / Phase 5t per-section counts (FermiHubbardDimer.lean) ---
\newcommand{\fhdTotal}{143}
\newcommand{\fhdWaveTwo}{13}
\newcommand{\fhdWaveThree}{6}
\newcommand{\fhdWaveFour}{22}
\newcommand{\fhdWaveFive}{22}
\newcommand{\fhdWaveSixABC}{38}
\newcommand{\fhdWaveSixExtra}{15}
\newcommand{\fhdWaveSeven}{19}
\newcommand{\fhdWaveEight}{8}

% \higherCurvatureThms and \higherCurvatureTests are auto-generated by
% scripts/update_counts.py from the canonical Lean source +
% pytest count and live in docs/counts.tex.

\begin{document}

\title{Higher-curvature structure from the Dirac heat kernel:
  microscopic predictions and observational ceilings}

\author{John G. Roehm}
\affiliation{Independent Researcher; \texttt{NetRxn137@gmail.com}}

\date{\today}

\begin{abstract}
We formalize, in the Lean~4 proof assistant, the basis change from the
3-scalar Wave~1 Dirac $a_4$ coefficients
$\{\Rsq,\, \Ric,\, \Riem\}$~\cite{ChristensenDuff1979,Vassilevich2003}
into Stelle's canonical
$\{\Rsq,\, \Csq,\, \GB\}$ basis~\cite{Stelle1977}, where
$\Csq$ is the Weyl-tensor-squared and $\GB$ the Gauss--Bonnet Euler
density (topological in $4$D~\cite{Lovelock1971}). The closed-form
coefficients $(\alpha, \beta, \gamma)$ at fermion count $\Nf$ are
fixed by a $3\times 3$ linear-system solve, with sign-definite
behaviour: $\alpha < 0,\ \beta < 0,\ \gamma > 0$ for any $\Nf > 0$.
The substantive cross-bridge to Wave~1 is encoded as a Lean
\texttt{ring}-closing identity. We then formalize a numerical
correctness-push that compares the predicted dimensionless
higher-curvature couplings (at SM-relevant counts $\Nf \in [24,27]$)
to the four canonical observational
ceilings~\cite{CalmetCapozzielloPryer2017,Berti2015}: LIGO/Virgo
speed-of-graviton, Hulse--Taylor binary-pulsar period decay,
Eöt-Wash short-range gravity, and the Cassini post-Newtonian solar-system
test. The Hulse--Taylor pulsar bound is the tightest at
$|\beta| \lesssim 10^{59}$; the predicted Wave~1 coefficients sit at
$\mathcal{O}(10^{-3})$, some $62$ orders of magnitude below. The
\texttt{HigherCurvatureStructure.lean} module ships
\higherCurvatureThms{} substantive theorems, zero \texttt{sorry}, and
zero new axioms; the Python mirror in \texttt{src/higher\_curvature/}
backs them with \higherCurvatureTests{} cross-checks.
\end{abstract}

\maketitle

\section{Introduction}
\label{sec:intro}

The Sakharov--Adler induced-gravity programme yields, at one loop
in the heat-kernel expansion, an Einstein--Hilbert action plus a
calculable tower of higher-curvature corrections. In Phase~6e Wave~1
we formalized the leading $a_2$ coefficient and its exact match to
the linearized 6a.1 emergent Newton constant at the Sakharov
baseline~$\alpha_{\mathrm{ADW}} = 1$~\cite{Roehm2026Wave1}. Wave~2
takes the next coefficient, $a_4$, and unpacks its three-scalar
expansion
\begin{equation}
  a_4(x) \;=\; \frac{\Nf}{(4\pi)^{2}}
    \Big[ c_R\, \Rsq + c_{\mathrm{Ric}}\, \Ric
      + c_{\mathrm{Riem}}\, \Riem \Big],
  \label{eq:a4}
\end{equation}
with the Christensen--Duff--Vassilevich rational
coefficients
$(c_R,\, c_{\mathrm{Ric}},\, c_{\mathrm{Riem}}) =
  \big(-\tfrac{5}{12 \cdot 180},\, \tfrac{7}{12 \cdot 180},\,
       -\tfrac{12}{12 \cdot 180}\big)$.
At the local effective-action level only \emph{two} combinations of
these scalars survive in $4$D: the trace-free Weyl-squared
$\Csq = \Riem - 2\Ric + \tfrac{1}{3}\Rsq$ and an arbitrary Ricci-trace
piece. The Gauss--Bonnet density
$\GB = \Rsq - 4\Ric + \Riem$ integrates to $32\pi^{2} \chi(M)$ on
closed $4$-manifolds and contributes nothing to the local equations
of motion.

The Wave~2 deliverable is twofold:
(i)~the closed-form basis change to Stelle's
$\{\Rsq, \Csq, \GB\}$ representation with sign-definite coefficients,
formalized as a \texttt{ring}-closing identity at the density level;
(ii)~the correctness-push comparison of the resulting dimensionless
couplings against current observational ceilings.

\section{Curvature basis: Gauss--Bonnet and Weyl-squared}
\label{sec:basis}

We formalize the two algebraically-independent curvature scalars
beyond $\Rsq$ as Lean definitions
\texttt{gaussBonnet4D} and \texttt{weylSquared4D}:
\begin{align}
  \GB &\;\equiv\; \Rsq - 4\Ric + \Riem ,
  \label{eq:gbdef} \\
  \Csq &\;\equiv\; \Riem - 2\Ric + \tfrac{1}{3}\Rsq.
  \label{eq:weyldef}
\end{align}
The Lean theorem
\texttt{weylSquared4D\_eq\_zero\_iff\_conformally\_flat}
encodes the conformal-flatness biconditional
$\Csq = 0 \iff \Riem = 2\Ric - \tfrac{1}{3}\Rsq$. Sanity check:
de~Sitter at Hubble parameter $H$ has
$(\Rsq,\Ric,\Riem) = (144 H^{4},\, 36 H^{4},\, 24 H^{4})$, giving
$\Csq = 24 - 72 + 48 = 0$ identically (de~Sitter is conformally flat).

The substantive algebraic engine for the Wave~2 basis change is the
identity
\begin{equation}
  \GB - \Csq \;=\; \tfrac{2}{3}\, \Rsq - 2\, \Ric ,
  \label{eq:gbweyl}
\end{equation}
formalized as
\texttt{gaussBonnet\_minus\_weyl\_eq\_R\_minus\_Ricci\_combination}
and closed by \texttt{ring} after unfolding both
Eqs.~\eqref{eq:gbdef}--\eqref{eq:weyldef}.

\section{Stelle-basis coefficients $(\alpha, \beta, \gamma)$}
\label{sec:coefs}

In the canonical $\{\Rsq, \Csq, \GB\}$ basis, the Wave~1 $a_4$ density
takes the form
\begin{equation}
  a_4(x) \;=\; \alpha(\Nf)\, \Rsq
    + \beta(\Nf)\, \Csq + \gamma(\Nf)\, \GB,
  \label{eq:a4stelle}
\end{equation}
where the coefficients $(\alpha, \beta, \gamma)$ are determined by
the linear system
\begin{equation}
  \begin{aligned}
    \alpha + \tfrac{\beta}{3} + \gamma &= c_R, \\
    -2\beta - 4\gamma &= c_{\mathrm{Ric}}, \\
    \beta + \gamma &= c_{\mathrm{Riem}}.
  \end{aligned}
  \label{eq:linsys}
\end{equation}
Solving Eq.~\eqref{eq:linsys} yields the closed forms
\begin{equation}
  \alpha = -\frac{\Nf}{324\,(4\pi)^{2}}, \quad
  \beta  = -\frac{41\, \Nf}{4320\,(4\pi)^{2}}, \quad
  \gamma = +\frac{17\, \Nf}{4320\,(4\pi)^{2}}.
  \label{eq:abg}
\end{equation}
Three sign-definite Lean theorems
(\texttt{a4\_alpha\_neg}, \texttt{a4\_beta\_neg},
\texttt{a4\_gamma\_pos}) encode the immediate consequence
\begin{equation}
  \Nf > 0 \;\Longrightarrow\;
    \alpha(\Nf) < 0,\ \beta(\Nf) < 0,\ \gamma(\Nf) > 0.
\end{equation}
The positive sign of the topological coefficient $\gamma$ is the
fingerprint of the chiral-anomaly contribution to the
Euler density: $\gamma > 0$ couples the Christensen--Duff Dirac
sector to the topological sector with a definite orientation.

\subsection*{Main basis-change identity}

The substantive cross-bridge between the original 3-scalar basis and
Stelle's $\{\Rsq, \Csq, \GB\}$ basis is the Lean theorem
\begin{verbatim}
theorem
a4_density_eq_a4_density_in_RC2GB_basis
  (N_f R_sq Ricci_sq Riemann_sq : ℝ) :
  a4_density N_f R_sq Ricci_sq Riemann_sq
    =
  a4_density_in_RC2GB_basis
    N_f R_sq Ricci_sq Riemann_sq
\end{verbatim}
proved by \texttt{ring} after unfolding the Wave~1 coefficients
\texttt{a4\_R\_sq\_coef}, \texttt{a4\_Ricci\_sq\_coef},
\texttt{a4\_Riemann\_sq\_coef}, the Wave~2 coefficients
\texttt{a4\_alpha}, \texttt{a4\_beta}, \texttt{a4\_gamma}, and the
definitions \texttt{weylSquared4D}, \texttt{gaussBonnet4D}. This is
drift-protection per the project's
\texttt{feedback\_python\_lean\_refs\_drift.md} memo: the proof body
literally invokes the Wave~1 names, so any future renaming would break
the build immediately.

\section{Observational ceilings on dimensionless higher-curvature
couplings}
\label{sec:bounds}

In the Stelle truncation
$L = (1/16\pi G_N)\big(R + \alpha \Rsq + \beta \Csq\big)$ the
dimensionless coefficients $(\alpha, \beta)$ are bounded by four
canonical observations~\cite{CalmetCapozzielloPryer2017,Berti2015}:
\begin{itemize}
  \item LIGO/Virgo speed-of-graviton (GW170817): $|\beta| \lesssim 10^{62}$
  \item Eöt-Wash short-range gravity (50~\textmu m): $|\alpha| \lesssim 10^{61}$
  \item Hulse--Taylor binary-pulsar period decay: $|\beta| \lesssim 10^{59}$
  \item Cassini post-Newtonian solar-system: $|\beta| \lesssim 10^{62}$
\end{itemize}
The Hulse--Taylor pulsar bound is tightest by $\sim 3$ orders of
magnitude. We formalize the four bounds as Lean definitions
\texttt{hc\_bound\_LIGO}, \texttt{hc\_bound\_SRG},
\texttt{hc\_bound\_pulsar}, \texttt{hc\_bound\_cassini}.

\section{Correctness-push: predictions vs ceilings}
\label{sec:cpush}

\begin{figure}
  \includegraphics[width=\columnwidth]{figures/fig_higher_curvature_obs_bounds.png}
  \caption{Higher-curvature predictions vs observational ceilings.
    \textbf{Left:} Magnitudes of the three Wave~1 Dirac $a_4$ Christensen--Duff
    coefficients (per $(4\pi)^{-2}$) at SM-relevant fermion counts
    $\Nf \in \{24, 27, 100\}$; the $\Riem$ coefficient dominates at every
    $\Nf$. \textbf{Right:} log-scale comparison to the four canonical
    observational ceilings; the Hulse--Taylor pulsar bound is tightest
    at $|\beta| \lesssim 10^{59}$, and the predicted $\mathcal{O}(10^{-3})$
    coefficient sits $\sim 62$ orders of magnitude below.
    Lean: \texttt{higher\_curvature\_below\_pulsar\_bound};
    \texttt{higher\_curvature\_predictions\_strictly\_positive};
    \texttt{H\_HigherCurvatureWithinObservationalBounds\_pulsar\_witness}.}
  \label{fig:obsbounds}
\end{figure}

The structural Wave~2 correctness-push is encoded as the Lean theorem
\begin{verbatim}
theorem higher_curvature_below_pulsar_bound
  {N_f : ℝ} (hN_pos : 0 < N_f) (hN_max : N_f ≤ 100) :
  |a4_R_sq_coef N_f|       < hc_bound_pulsar ∧
  |a4_Ricci_sq_coef N_f|   < hc_bound_pulsar ∧
  |a4_Riemann_sq_coef N_f| < hc_bound_pulsar
\end{verbatim}
proved by combining (i)~the bound $\fbox{1 < (4\pi)^2}$ from
\texttt{Real.pi\_gt\_three} and \texttt{nlinarith} (helper lemmas
\texttt{fourPiSq\_gt\_one} and \texttt{fourPiSqInv\_lt\_one});
(ii)~the trivial numerical estimate
$100 \cdot 12 / (12 \cdot 180) \cdot 1 < 1$; and
(iii)~the chain $1 < 10^{59}$.
The proof body invokes the Wave~1 \texttt{a4\_R\_sq\_coef},
\texttt{a4\_Ricci\_sq\_coef}, \texttt{a4\_Riemann\_sq\_coef} by name
(P6 cross-module bridge integrity).

The 3-conjunct bundle is \emph{not} a P2-redundancy: each conjunct
invokes a distinct Wave~1 coefficient and a distinct positivity proof
(the $R^2$ and $\Riem$ coefficients are negative and need an
explicit \texttt{abs\_neg} step; the $\Ric$ coefficient is positive
and uses \texttt{abs\_of\_pos} directly).

\paragraph*{Anchoring the "62 orders below" claim.}
At the SM-with-$\nu_R$ count $\Nf = 27$ the dominant Riemann${}^2$
coefficient evaluates to
\begin{equation}
  |c_{\mathrm{Riem}}(\Nf=27)|
  \;=\; \frac{27}{180\,(4\pi)^{2}}
  \;\approx\; 9.49\times 10^{-4},
  \label{eq:c-riem-numeric}
\end{equation}
and the same form bounds the smaller $|c_R|, |c_{\mathrm{Ric}}|$
in the same dimensionless natural-units convention.  Comparing
against the binary-pulsar ceiling $|\beta| \lesssim 10^{59}$ gives
$10^{59}/9.49\times 10^{-4} \approx 10^{62}$, i.e.\ the predicted
coupling sits ${\sim}62$ orders of magnitude below the tightest
observational bound.  The numerical evaluation is reproduced in
\texttt{tests/test\_higher\_curvature.py} (\texttt{TestObservationalBounds}).

A complementary falsifier
\texttt{higher\_curvature\_predictions\_strictly\_positive} ensures
that the bounds are \emph{not} passed by trivial vanishing: every
Wave~1 $a_4$ coefficient is non-zero for $\Nf > 0$. This rules out
the unhelpful reading "all bounds are passed because all predictions
are zero".

The Lean Prop predicate
\texttt{H\_HigherCurvatureWithinObservationalBounds B} is parameterised
by an upper bound $B$:
\[
  H_{\mathrm{HC}}(B) :\;
    0 < B \;\wedge\; \forall\, \Nf,\ 0 < \Nf \le 100,
\]
\[
  \quad\quad |a_4(R^2)| \le B,\ |a_4(\Ric)| \le B,\ |a_4(\Riem)| \le B.
\]
Unlike a typical tracked-hypothesis Prop, this predicate is
\emph{fully discharged} for $B = 10^{59}$ inside the module: the
Lean theorem
\texttt{H\_HigherCurvatureWithinObservationalBounds\_pulsar\_witness}
proves $H_{\mathrm{HC}}(10^{59})$ by invoking the correctness-push,
so consumers receive a closed proof rather than an undischarged
hypothesis. The parameterisation over $B$ is a forward-compatibility
hook: future tighter ceilings (e.g.~from next-generation pulsar
timing) can be discharged at the new $B$ without changing downstream
consumers.

\section{Lean module structure}
\label{sec:lean}

The module
\texttt{lean/SKEFTHawking/HigherCurvatureStructure.lean} ships
\higherCurvatureThms{} substantive theorems organised in six sections:
\begin{description}
  \item[§1] \texttt{gaussBonnet4D}, \texttt{weylSquared4D}
    definitions plus the conformal-flatness biconditional and the
    \texttt{gaussBonnet\_minus\_weyl} algebraic identity (2 theorems).
  \item[§2] Stelle-basis coefficients \texttt{a4\_alpha},
    \texttt{a4\_beta}, \texttt{a4\_gamma} (definitions) plus
    sign-definiteness for $\Nf > 0$ (3 theorems).
  \item[§3] Density-level definitions \texttt{a4\_density},
    \texttt{a4\_density\_in\_RC2GB\_basis}, plus the substantive
    basis-change identity \texttt{a4\_density\_eq\_a4\_density\_in\_RC2GB\_basis}
    (1 theorem).
  \item[§4] Helper bounds \texttt{fourPiSq\_gt\_one},
    \texttt{fourPiSqInv\_lt\_one} (2 theorems).
  \item[§5] Observational-ceiling \texttt{def}s.
  \item[§6] Correctness-push \texttt{higher\_curvature\_below\_pulsar\_bound},
    falsifier \texttt{higher\_curvature\_predictions\_strictly\_positive},
    plus the tracked-Prop witness
    \texttt{H\_HigherCurvatureWithinObservationalBounds\_pulsar\_witness}
    (3 theorems).
\end{description}

The module ships zero \texttt{sorry} and zero new axioms beyond the
Lean~4 kernel's standard \texttt{propext}, \texttt{Classical.choice},
\texttt{Quot.sound}. The Python mirror in
\texttt{src/higher\_curvature/} provides numerical evaluation of the
$\{\Rsq, \Csq, \GB\}$ basis, the Stelle coefficients, and the
observational-bound check; a \higherCurvatureTests{}-case
\texttt{pytest} suite cross-checks each Lean theorem against the
Wave~1 numerics.

\section{Conclusion}
\label{sec:conclusion}

The Wave~1 Christensen--Duff Dirac $a_4$ coefficient triple is
re-expressed in Stelle's canonical $\{\Rsq, \Csq, \GB\}$ basis with
closed-form sign-definite Stelle coefficients
$(\alpha, \beta, \gamma)$. The basis change is recorded as a Lean
\texttt{ring} identity at the density level (cross-bridge to Wave~1).
A correctness-push theorem certifies that, in the natural fermion-count
window $0 < \Nf \le 100$, every predicted dimensionless higher-curvature
coefficient sits $\sim 62$ orders of magnitude below the tightest
observational ceiling (Hulse--Taylor binary pulsar) --- no tension
with current observation. The complementary falsifier rules out the
trivial reading.

Wave~3 (\texttt{NonlinearDiffInvariance.lean}) consumes the
$\{\Rsq, \Csq, \GB\}$ basis to formalize order-by-order
diff-invariance of the heat-kernel effective action; the structural
correctness-push there is whether $a_4$ contributions to the
nonlinear effective Einstein equations are diff-invariant or
not~--- a clean falsifiable test of the ADW emergent-gravity claim.

\begin{thebibliography}{99}

\bibitem{ChristensenDuff1979}
S.~M.~Christensen and M.~J.~Duff,
\emph{New gravitational index theorems and supertheorems},
Nucl. Phys. B \textbf{154}, 301 (1979).

\bibitem{Vassilevich2003}
D.~V.~Vassilevich,
\emph{Heat kernel expansion: user's manual},
Phys. Rep. \textbf{388}, 279 (2003)
[arXiv:hep-th/0306138].

\bibitem{Stelle1977}
K.~S.~Stelle,
\emph{Renormalization of higher-derivative quantum gravity},
Phys. Rev. D \textbf{16}, 953 (1977).

\bibitem{Lovelock1971}
D.~Lovelock,
\emph{The Einstein tensor and its generalizations},
J. Math. Phys. \textbf{12}, 498 (1971).

\bibitem{CalmetCapozzielloPryer2017}
X.~Calmet, S.~Capozziello and D.~Pryer,
\emph{Gravitational effective action at second order in curvature
and gravitational waves},
Eur. Phys. J. C \textbf{77}, 589 (2017)
[arXiv:1708.08253].

\bibitem{Berti2015}
E.~Berti \emph{et al.},
\emph{Testing general relativity with present and future astrophysical
observations},
Class. Quantum Grav. \textbf{32}, 243001 (2015)
[arXiv:1501.07274].

\bibitem{Roehm2026Wave1}
J.~G.~Roehm,
\emph{Formal heat-kernel calibration of induced Newton's constant from
the ADW Dirac substrate}
(in preparation, Phase 6e Wave 1).

\end{thebibliography}

\end{document}
